\section{What goes into the model}

Section 1.4 introduced the three workbooks that define a scenario.
In the 2023 dataset, four climate scenarios, which differ in their projections and calibration, can be combined with three geopolitical scenarios, which differ in whether Russian gas reaches Europe.

The inputs are not all of the same kind.
Table \ref{tab:ggm-inputs} groups the numeric values in one 2023 scenario by where they come from, drawing on the documentation and on a recorded working session with one of the model's developers.\footnote{Workflow session with Lukas Barner (TU Berlin), 23 October 2026, in which he walked through the three workbooks and the calibration on his own machine.}

\begin{table}[htbp]
\centering
\small
\caption{Kinds of input in one scenario of the 2023 GGM dataset (our count).}
\label{tab:ggm-inputs}
\begin{tabularx}{\textwidth}{@{}lrX@{}}
\toprule
Kind of input & Values & Where it comes from \\
\midrule
Physical facts & 15,200 & Capacities, lengths and distances from industry and agency databases \\
Reference quantities & 2,450 & Production and consumption per node and year: a storyline aligned with a published outlook, adjusted node by node \\
Demand shares & 6,950 & Sector and seasonal shares, from earlier studies and analogies between countries \\
General parameters & 26 & Elasticities, base costs, loss and discount rates \\
Scenario settings & 1,900 & Expansion limits, market power, sales restrictions \\
Calibrated values & 7,050 & Set by calibration; about 2,630 differ from their neutral value \\
\bottomrule
\end{tabularx}
\end{table}

Only the first kind is collected in the ordinary sense.
Pipeline and storage capacities come from ENTSOG and GIE for Europe, from the EIA for the United States and from the Global Energy Monitor for the rest of the world, and are mapped from countries or states to model nodes; shipping distances are computed from terminal coordinates.
This work is laborious, because sources report in their own units and geographies, but it follows rules, and the developers already use AI tools for the mapping.

The reference quantities are different.
They are anchored to the aggregates of an outlook such as the IEA's World Energy Outlook, but the analysts adjust them wherever they find the outlook implausible, and split country totals across nodes using regional statistics or, where none exist, proxies such as regional GDP.
The result cannot be reverse-engineered from the outlook, and is not meant to be: it is the storyline of the scenario, and in the developers' words there is no right or wrong version of it.
Elasticities and the sector and seasonal shares are similar, since no data set of them exists; they are carried over from earlier work or rebuilt for each study.
Preparing the projections for a new study takes an experienced analyst about a week.

Part of the source data is proprietary, so the model is published without it, and the documentation describes the processing in enough detail for others to rebuild the data \citep[p.~8]{egging2019global}, a common compromise in energy modelling \citep{pfenninger2017importance}.
The scenario settings express which world is being modelled, and the calibrated values are the subject of Section 3.5.

% TODO: Agree with Franziska how to cite the workflow session (footnote as personal communication, or an appendix with the session notes).
